% ----------------------------------------------------------
% Seção 6 — Ferramentas e fluxo de trabalho
% ----------------------------------------------------------
\chapter{Ferramentas e Fluxo de Trabalho}\label{cap:ferramentas}

Escrever em \LaTeX\ é parecido com programar: o texto-fonte é compilado para gerar o PDF. Práticas consolidadas no desenvolvimento de software, como controle de versões e automação de tarefas, também se aplicam à escrita acadêmica \cite{wilson2017}.

\section{Compilação}\label{sec:compilacao}

O modelo é compilado com o \texttt{latexmk} \cite{latexmk}, que executa o \texttt{pdflatex} e o programa de bibliografia quantas vezes forem necessárias para resolver citações, referências cruzadas e sumário. Na pasta do modelo:

\begin{minted}{bash}
latexmk main.tex      # compila e gera main.pdf
latexmk -C            # apaga os arquivos gerados
\end{minted}

O programa de bibliografia depende da versão da ABNT escolhida: o Bib\TeX\ para a versão 2002, com o pacote \texttt{abntex2cite} \cite{abntex2cite}, e o Biber para as versões 2018 e 2023, com os pacotes \texttt{biblatex} \cite{biblatex} e \texttt{biblatex-abnt} \cite{biblatexabnt}. O \texttt{latexmk} detecta qual deles usar, mas, ao trocar de versão, apague os arquivos gerados com \texttt{latexmk -C} antes de compilar.

O pacote \texttt{minted}, usado para código-fonte, precisa do Python 3 instalado. Se não for usar código-fonte, remova a linha \verb|\usepackage{minted}| e os exemplos que usam o ambiente \texttt{minted}.

\subsection{Overleaf}

O Overleaf é um editor de \LaTeX\ on-line. Para usar o modelo nele, baixe o arquivo \texttt{.zip} da versão da ABNT desejada na página de \textit{releases} do repositório e use a opção \textit{New Project} > \textit{Upload Project}.

\subsection{Compilação de todos os modelos}

No repositório do projeto, o \texttt{Makefile} compila todos os modelos em todas as versões da ABNT e gera os pacotes de distribuição:

\begin{minted}{bash}
make pdf ABNT=2002   # todos os modelos na versão escolhida
make todas           # todos os modelos nas três versões
make dist            # pacotes .zip e PDFs em dist/
\end{minted}

\section{Gerenciadores de referências}\label{sec:gerenciadores}

Gerenciadores de referências organizam a bibliografia da pesquisa e exportam o arquivo \texttt{.bib}. O JabRef\footnote{Disponível em: \url{https://www.jabref.org/}.} edita arquivos \texttt{.bib} diretamente. O Zotero\footnote{Disponível em: \url{https://www.zotero.org/}.}, com a extensão Better Bib\TeX, mantém o arquivo \texttt{.bib} atualizado automaticamente.

\section{Controle de versões}\label{sec:git}

O Git registra o histórico do trabalho e permite recuperar qualquer versão anterior do texto \cite{chacon2014}. Um fluxo mínimo é:

\begin{minted}{bash}
git init                              # uma vez, na pasta do trabalho
git add -A
git commit -m "Escreve a seção de metodologia"
\end{minted}

Hospedar o repositório em um serviço como o GitHub cria uma cópia de segurança e facilita a revisão pelo orientador. Registrar também os dados, os scripts de análise e as versões dos programas usados permite que os resultados do trabalho sejam reproduzidos \cite{sandve2013}.
